\chapter{Introdução}
\label{cap:introducao}
% OKANAGAN - PRA ENSINO SUPERIOR
% DASU - POLITICA DE SAUDE MENTAL
% PROMOÇAO DE SAUDE E 

% DOCUMENTO OSM - NA CONTEXTUALIZAÇAO PROMOÇAO A SAUDE MENTAL
% POS PANDEMIA

% vigilancia

% ------------

% Anotações e topicos para introdução, NÃO APAGAR TO TRABALHANDO NISSO:

% O regimento da ReBraUPS descreve a Universidade Promotora de Saúde UPS como uma rede de instituições de ensino superior que visa promover a saúde e criar ambientes universitários saudáveis.

% A UPS combate práticas que comprometam o bem-estar e a qualidade de vida, promovendo ações que fortalecam ambientes saudáveis.

% Tendo isso em mente, a ReBraUPS segue os principios da Carta de Okangan (2015), que fala que a promoção da saúde é identificada como ações cotidianas no âmbito universitário, dando poder a abordagens envolvendo vozes da comunidade acadêmica.

% Além disso, segundo a Carta de Edmonton (2005) define que as universidades promotoras de saúde possuem responsabilidade para contribuir com o aumento da saúde e do bem estar, eles lideram essa parada tlgd? e a ReBraUPS busca fortalecer o papel da universidade para combater a falta de saúde, sendo assim, a Carta de Edmonton define que as universidades sejam agentes de mudança na promoçao a saúde através de diversas ações.


% Aprendizagem - psicologas escolares e pedagogas
% Ansiedade, sofrimento - psicologas clinicas

% Contrução de sentido de pertencimento - planejamento de estudo...

% Universidade Promotora da Saúde (UPS) é uma abordagem institucional e estratégica que busca incorporar promoção da saúde no ensino, pesquisa e extensão dentro da universidade. Consideramos uma universidade como UPS quando ela adota um conjunto de estragias que visam promover ambientes saudáveis e sustentaveis.

% Rede Brasileira de Universidades Promotoras de Saúde (ReBraUPS) é um conjunto de instituições que se comprometem a se organizar para atuar como UPS.

% OKONAGAN, 2015 - o ensino superior exerce um grande papel no desenvolvimento de indivíduos, comunidades, sociedades e culturas, tendo a oportunidade e a responsabilidade de oferecer uma educação que transforme. Podemos falar sobre influenciar os alunos adquirirem hábitos saudaveis para o resto da vida...

O ministro canadense Marc Lalonde, em 1974, criou o primeiro documento a adotar a denominação de promoção da saúde, em um informe conhecido como "O Informe Lalonde". O mesmo destaca a importância da promoção da saúde no meio universitário \cite{rebraups2022}. O autor estimula a ideia de que a saúde está diretamente relacionada com estilo de vida, organização assistencial, meio ambiente e biologia humana, sendo assim, o conceito de saúde vai além do sistema tradicional, englobando diversos aspectos extras \cite{rezende_promocao_2022}.

Além disso, no relatório de Lalonde é adotado pela primeira vez o termo Determinantes Sociais de Saúde (DSS), que são fatores que determinam a saúde de uma pessoa, como boa alimentação, possibilidade de lazer, boas condições de trabalho, hábitos saudáveis, entre outros \cite{heidmann_promocao_2006}. 

Indubitavelmente, o ensino superior tem um papel fundamental no desenvolvimento de indivíduos, sociedade e culturas, o mesmo possui a responsabilidade de oferecer uma educação que transforma a vida dos estudantes, como descrito na Carta de Okanagan. Sendo assim, é essencial que as Instituições de Ensino Superior (IES) estimulem a promoção da saúde, dessa forma foi criado o movimento de Universidades Promotoras de Saúde (UPS), cujo intuito é reunir instituições que tenham o objetivo de promover a saúde e debater soluções e avanços \cite{rezende_promocao_2022}. 

Tendo isso em mente, segundo \cite{schleich_integracao_nodate}, em "Integração na Educação Superior e Satisfação Acadêmica de Estudantes Ingressantes e Concluintes", define importantes pontos que as IES devem adotar como melhores meios de inovação tecnológica, novos espaços educativos e maior conhecimento sobre os estudantes, sendo esse último um dos fundamentos para a criação desse trabalho.

Sendo assim, como descrito por \citeonline{schueller_2019}, a utilização de prática de \textit{mindfulness}, onde o paciente foca a atenção no momento presente e diários de humor, apresenta resultados positivos, auxiliando na redução do estresse e promovendo o bem-estar emocional. Nesse contexto, é possível a utilização de meios tecnológicos para auxiliar a promoção da saúde dentro das IES.

% MINHA INTENÇÃO:

% Introdução: Explicar que é dever da faculdade a criar meios pra promover a saúde
% Contextualização: mostrar que a faculdade é um ambiente propicio a adquirir problemas de saúde 
% Justificativa: juntar os dois provando que há a necessidade de criação de um aplicativo pra auxiliar a adoção de hábitos saúdaveis e auxiliar profissionais de saúde.


\section{Contextualização}


O Brasil, segundo o Instituto Brasileiro de Geografia e Estatística (IBGE), em 2014, a estimativa de depressão entre indivíduos entre 18 a 29 anos corresponde a 3,9\% da população, correspondendo a um grupo de jovens adultos. Boa parte disso se deve à necessidade de tomada de decisões, pressão social, entre outros fatores \cite{souza_condicoes_2017}.

Esse período é conhecido por possuir um dos maiores riscos psicológicos, podendo gerar gatilhos emocionais e estresse duradouro \cite{silveira_saumental_2011}. Também citado por \citeonline{bolsoni-silva_o_2014}, a universidade é capaz de agravar problemas de saúde pré-existentes ou aumentar a probabilidade do seu surgimento. Sendo assim, a universidade é um ambiente estressante, principalmente no primeiro semestre. A solidão, um fator crucial para depressão, ocorre principalmente no primeiro ano da universidade devido a um novo ambiente e à falta de experiências anteriores em contextos desafiadores e exigentes \cite{rhodes_loneliness_2014}.

Ressalta-se a importância dos hábitos saudáveis, lazer, presença religiosa e atividade física, que ajudam a promover o bem-estar e são possíveis formas de lidar com estados emocionais negativos. Outro fator crucial para o desenvolvimento de problemas psicológicos é a falta de hábitos saudáveis e menor qualidade de vida nos domínios psicológico e ambiental \cite{barroso_solidao_2019}.

Além disso, \citeonline{barroso_solidao_2019} destaca como os hábitos de vida influenciam diretamente a relação com a solidão e a depressão. Comportamentos como acordar cansado e realizar menos refeições ao longo do dia, entre outros, revelaram um impacto significativo na qualidade de vida dos participantes.

Outro aspecto que influencia a saúde emocional é o suporte social que, como defendido por \citeonline{siqueira_construcao_2008} em "Construção e validação da escala de percepção de suporte social", é definido como “o processo no qual alguém pode conseguir ajuda de ordem instrumental, financeira ou emocional”. Esse apoio é considerado essencial para a promoção e manutenção da saúde mental.

\section{Justificativa}


% DOCUMENTO OSM - NA CONTEXTUALIZAÇAO PROMOÇAO A SAUDE MENTAL
% FALAR DA VIGILANCIA NA CONTEXTUALIZAÇÃO
%desatualizado fonaprace é precario usar outra referencia


% De acordo com o Fórum Nacional de Pró-reitores de Assuntos Comunitários e Estudantis (FONAPRACE) em 2019, o impacto da solidão e depressão e percepção ruim do curso e relacionamento com colegas de sala aumentam o risco de desistência, reprovações e trancamentos na universidade.

Devido aos altos índices de risco psicológico destacados por \citeonline{silveira_saumental_2011}, o ambiente universitário revela-se potencialmente hostil à saúde mental dos estudantes, exigindo a implementação de estratégias de promoção à saúde que vão além de ações institucionais, como palestras, e incluam também abordagens personalizadas e voltadas ao indivíduo, como diários de humor \cite{schueller_2019}.

Nesse contexto, a vigilância participativa configura-se como uma abordagem essencial no monitoramento da população acadêmica, ao permitir a identificação precoce e a resposta eficiente a situações que afetam coletivamente a saúde mental \cite{lealneto2016digital}. Sendo assim, a fim de fornecer uma maneira de auxiliar a promoção de saúde nas IES, é fundamental a utilização deste conceito.

\vspace{23em}

Além disso, a adoção de hábitos saudáveis mostra-se eficaz no enfrentamento e na mitigação dos impactos das doenças mentais, conforme apontado por \citeonline{barroso_solidao_2019}. Ou seja, A promoção desses hábitos pode refletir positivamente no desempenho acadêmico e na permanência dos estudantes nos cursos. Além disso, com base no conceito de UPS se da a necessidade de criar meios para auxiliar a promoção a saúde nas IES, como mencionado por \citeonline{rezende_promocao_2022}, não há uma fórmula única para promover a saúde no ambiente universitário.

Dessa forma, é fundamental a criação de uma ferramenta para auxiliar a adoção de hábitos saúdaveis e promover a saúde mental no meio acadêmico e uma ferramenta para auxiliar profissionais de saúde acadêmicos a identificar alunos em risco. Ademais, tendo em vista o conceito de suporte social é possível a utilização da tecnologia para criação de um software que realize esse papel \cite{siqueira_construcao_2008}. 



% Sendo assim, com base no conceito de ups se da a necessidade de criar meios para auxiliar a ps na faculdade
% Ademais segundo... Se da necessidade da criaçao de uma aplicaçao para auxiliar a triagem de alunos

\section{Questão de Pesquisa}

  
  Para fundamentar o desenvolvimento deste trabalho, a questão de pesquisa foi formalmente definida: Como promover intervenções personalizadas na saúde mental de estudantes universitários por meio de um aplicativo mobile baseado na análise de dados sobre hábitos de vida e estados emocionais?




\section{Objetivos}
Este TCC possui um objetivo geral e outros oito objetivos específicos que serão aprofundados nesta seção.

\subsection{Objetivo Geral}

Este trabalho possui como objetivo geral o desenvolvimento de um aplicativo mobile que visa a promoção da saúde entre os membros da comunidade universitária, por meio da coleta e análise de dados sobre seus hábitos de vida e estados de saúde mental.

\subsection{Objetivos Específicos}

São objetivos específicos deste trabalho:

\begin{itemize}

\item Contribuir para a identificação de obstáculos no bem-estar dos indivíduos de determinada instituição; 

\item Coletar dados relevantes que subsidiem a elaboração de projetos sociais e de saúde pela instituição;

\item Disponibilizar gráficos de progressos individuais para os usuários do aplicativo;

\item Oferecer orientações personalizadas de acordo com o progresso individual para melhoria de aspectos que apresentem evolução insatisfatória;

\item Disponibilizar um recurso passível de ser utilizado de forma ampla pela comunidade acadêmica para a aplicação de questionários relacionados à saúde da comunidade universitária;

\item Servir como uma ferramenta para disseminação de informações relevantes sobre o tema de saúde entre os indivíduos do campo universitário;

\item Propor uma estratégia de Promoção da Saúde aos estudantes e funcionários universitários;

\item Criar indicadores institucionais que avaliem o impacto de ações de promoção da saúde implementadas na comunidade universitária.
\end{itemize}

% \section{Metodologia}
% % A metodologia deste TCC é dividida em metodologia de pesquisa e metodologia de desenvolvimento, que serão descritas brevemente a seguir.

% \subsection{Metodologia de Pesquisa}

% Conforme descrito por \citeonline{zanella2006metodologia}, a pesquisa apresenta diversos tipos relacionado à sua natureza, objetivos, abordagem e procedimentos. Ainda de acordo com o material, é possível dizer que a presente pesquisa possui:

% \begin{itemize}

% \item \textbf{Natureza:} aplicada, por ter como objetivo a solução de um problema humano;
% \item \textbf{Objetivos:} exploratória, em razão de ser uma pesquisa que busca ampliar o conhecimento sobre fenômenos da comunidade acadêmica, atrás de conhecer suas necessidades e buscar tratá-las; 
% \item \textbf{Abordagem:} quantitativa, visto que caracteriza-se pela ampla utilização de instrumental estatístico fundamental na análise de dados coletados;
% \item \textbf{Procedimentos:} procedimento bibliográfico para investigação do referencial teórico aliado a produção tecnológica no contexto de desenvolvimento de software.

% \end{itemize}

% \subsection{Metodologia de Desenvolvimento}

% A metodologia de desenvolvimento deste trabalho segue princípios de práticas ágeis e une características do Extreme Programming (XP) ao Scrum, Kanban e Lean Inception. 

% No que diz respeito a metodologias ágeis, os atributos de programação em pares, aplicação de testes automatizados com TDD e a refatoração foram aproveitados da metodologia XP \cite{highsmith_2000}. Já as características extraídas da metodologia Scrum foram os conceitos de sprints iterativas e incrementais, junto das reuniões diárias \cite{schwaber_scrum_2013}. Vale destacar que esta última metodologia coube adaptações no contexto de uma equipe reduzida. 

% O Kanban, por sua vez, foi adotado como uma ferramenta de gestão visual para acompanhar o andamento das tarefas. Por fim, o Lean Inception estruturou a definição e priorização do Mínimo Produto Viável (MVP), que possibilitou a conformidade entre os envolvidos sobre a ideia do produto final.

% Maiores detalhes sobre a metodologia de desenvolvimento serão fornecidas no Capítulo 3 deste projeto.

\section{Organização do Trabalho}
O presente Trabalho de Conclusão de Curso (TCC) está dividido nos seguintes capítulos:

\begin{itemize}

\item \textbf{Capítulo \textcolor{blue}{1} - Introdução:} o capítulo introdutório apresenta o cenário de aplicação da proposta, a justificativa da realização da pesquisa, a questão que norteia este trabalho, os objetivos pretendidos e as metodologias empregadas no desenvolvimento deste projeto.

\item \textbf{Capítulo \textcolor{blue}{2} - Referencial Teórico:} neste capítulo são expostos os conceitos principais que cercam a temática deste trabalho tratados em seis tópicos principais, que abordam definições sobre saúde, hábitos de vida, aplicativos já existentes, escala de likert e demais conceitos relacionados ao tema.

\item \textbf{Capítulo \textcolor{blue}{3} - Metodologia:} capítulo voltado para a apresentação da metodologia de pesquisa e de desenvolvimento do projeto, em que são detalhadas as abordagens utilizadas para a implementação eficaz da ferramenta de software.

\item \textbf{Capítulo \textcolor{blue}{4} - Desenvolvimento:} este capítulo é destinado à apresentação dos detalhes do desenvolvimento do software e os ajustes realizados para garantir o êxito da implementação. 

\item \textbf{Capítulo \textcolor{blue}{5} - Considerações Finais:} este capítulo final expõe os resultados obtidos no trabalho até então, assim como as sugestões para desenvolvimentos futuros. 

\end{itemize}